% =====================================================================
% Abstract for EMMET-Newt paper (v2.0+).
% Target venue: IEEE TNSM, Computer Networks (Elsevier).
% Length: 200-280 words.
% Tone: framework contribution, mechanistic vocabulary, no hype.
% =====================================================================

\begin{abstract}
Adaptive routing algorithms react to network congestion through edge
state shared across packets, but lack two ingredients we find
empirically useful: \emph{per-packet history} (each packet's path so
far) and \emph{persistent loss memory} (which edges recently dropped
packets, beyond instantaneous load). We introduce \textbf{EMMET-Newt},
a framework that builds two complementary mechanisms on a gradient
router. The first, \emph{mass-aware routing}, gives each packet an
internal scalar (its \emph{mass}) that grows with traversed congestion
and scales the cost of candidate edges, biasing more-traveled packets
toward less-loaded paths. The second, the \emph{reaction-memory}
mechanism, increments a persistent per-edge bloodstain on every packet
drop so subsequent packets evade recently-failing edges---a discrete
Newton-III feedback on the routing potential. We instantiate the
framework as a constrained dynamic-programming router and a bursty
variant with delayed feedback (releasing accumulated memory at burst
boundaries, the conservative reading).
On 600 paired scenarios under Poisson workload (Erd\H{o}s--R\'enyi, scale-free,
small-world, regular lattice and real WAN backbones), the mass-aware
router cuts congestion losses against the Dijkstra-with-potential
baseline by 43.5\% on G\'EANT (95\% CI $[28.4, 57.2]$\%), 61.5\% on
Watts--Strogatz (95\% CI $[36.9, 79.0]$\%) and 93.0\% on a regular
grid (95\% CI $[83.7, 100]$\%). On bursty cross-topology runs ($n=30$
per cell), adding reaction memory improves the mass-aware router
further and outperforms CONGA-WAN on structured topologies (+61\% on
a $10\times 10$ grid with the conservative DELAY variant), while
remaining competitive but not dominant on Abilene and G\'EANT.
A temporal-ripple extension inspired by wave propagation closes the
gap on GÉANT to a statistical tie ($-3.3\%$, CI $[-10.1, +3.4]$\%)
and extends the Grid$10\times 10$ headline to $+66.7\%$
($[+56.9, +74.4]$\%).
We extend to $n=1{,}000$ and document the saturation regime where
both arms reach 100\% delivery and no signal remains. The framework
also accommodates two theoretical siblings (collective pressure,
pairwise repulsion) that we sketch but do not claim empirically.
\end{abstract}

\begin{keyword}
adaptive routing \sep
load-aware routing \sep
congestion control \sep
stateful packet processing \sep
reaction memory \sep
network simulation
\end{keyword}
